% ===== 问题四正文 =====
% 问题四正文；六张示意图的 TikZ 宏在 figures/q4_schematics.tex 中定义。
% 算法依据 Q4_V6_Lite 简化版；实验数字只引用该版本的本地实测记录。
\subsection{问题四：全向与定向混合干扰源的搜索清除策略}\label{sec:q4}

问题四和问题三的任务一样，都是从原点出发把 $10$ 到 $16$ 个源全部找到并清除，区别只在源的种类：一部分源是定向的，只朝发射方向两侧各 $90^\circ$ 的范围辐射信号，而且哪些源是定向的、朝哪个方向发射都不知道。这一个改动让第三问里三处依赖“全向接收”的做法都不再成立：一是无信号不再说明离得远，第三问用负观测缩小定位区域和省略检测的办法不能再用；二是覆盖站的布置只考虑了距离，对背对着检测点发射的源不起作用；三是贴着边界朝外发射的源，在场地里面任何位置都收不到，只能从场地外侧去听。本节先把定向源的信号模型和这几个后果交代清楚，再给出能同时对付两类源的发现覆盖条件、定位区域的更新规则和任务安排方法，然后说明求解流程，最后用本地演练数据检验效果。第三问已经建立、在本问仍然成立的部分（示向度楔形、包围圆、可靠清除判据、滚动规划的框架）这里只作引用，不再重复。

\subsubsection{模型准备}\label{sec:q4-1}

\textbf{（1）定向源的信号模型。}记某个频道的源位于 $p$，接收半径为 $R_f$，它是 $1000$ 到 $1500$ m 之间的一个固定值，但不知道具体是多少。全向源在距离不超过 $R_f$ 的任何位置都能被收到。定向源多一个未知的发射方向，记为单位向量 $u$；它只向 $u$ 两侧各 $90^\circ$ 的范围辐射，也就是只向以“过 $p$ 且垂直于 $u$ 的直线”为边界、包含 $u$ 那一侧的半平面辐射。于是在位置 $x$ 能收到定向源信号的条件是

\begin{equation}\label{eq:q4-1}
\|x-p\|\le R_f\quad\text{且}\quad u\cdot(x-p)\ge0 .
\end{equation}

第二个条件说的就是 $x$ 落在那个半平面里，图\ref{fig:q4-emission}(a) 画出了它的样子。收到信号后返回的示向度、$5$ m 内的近场反馈，以及 $20$ m 内才能完成的光学定位与清除，都和前面几问相同；光学清除不受发射方向影响，只要走到源的 $20$ m 以内就能清除。

\textbf{（2）无信号的两种原因。}对全向源，某处收不到信号只有一个解释：离源超过了 $R_f$。第三问正是靠这一点，把无信号点和有信号点配起来得到距离不等式\eqref{eq:q3-3}，又用“比某个无信号点离源更远”的推断省掉了一批检测（式\eqref{eq:q3-14}）。对定向源，无信号还可能是因为检测点在发射范围之外，图\ref{fig:q4-emission}(a) 中的两个无信号点就是这两种情况：一个离得近但在源的背面，一个在正面但离得远。既然分不清是哪种原因，一次无信号本身对源的位置就没有任何约束，第三问的这两条推断都不能沿用。有信号的观测则和以前完全一样：示向度给出 $\pm1^\circ$ 的楔形，收到信号说明距离不超过 $1500$ m，所以第二问建立的定位区域和它的更新方法在本问照常使用。

\textbf{（3）贴着边界朝外发射的源。}把式\eqref{eq:q4-1} 用在场地边界上，就能看到本问最棘手的情形。若源在半径 $1800$ m 的圆周上，$u$ 正好指向圆外，那么发射范围的边界就是圆在该点的切线，整个场地都在切线的另一侧，场内任何位置都收不到它，见图\ref{fig:q4-emission}(b)。要发现这样的源，检测点必须放到场地外侧去。题目只说干扰源都在目标区域内，没有限制机器狗的活动范围，本文因此允许机器狗在场地边界外不远处停下检测；后面的检测站布局有一环放在半径 $1865$ m 的圆上，就是为了这类源。

\begin{center}
\begin{minipage}{\linewidth}
\centering
\QFourEmissionFigure
\captionof{figure}{定向源的信号模型。(a) 定向源只向半平面 $u\cdot(x-p)\ge0$ 辐射，无信号既可能因为超出接收距离，也可能因为落在发射范围之外。(b) 贴着边界朝外发射的定向源，场内检测点全部无信号，只有场外的检测点能收到。}
\label{fig:q4-emission}
\end{minipage}
\end{center}

\textbf{（4）沿用的记号与信息状态。}目标圆域 $\Omega$、频道 $f$、当前位置 $x_t$、定位区域的外包多边形 $\widehat C_f$、包围圆半径 $\rho(\widehat C_f)$ 和圆心 $z(\widehat C_f)$ 都沿用第二、三问的记号；总耗时仍按式\eqref{eq:q3-2} 由移动、换频、检测、光学定位和清除五项累加，评价指标是清除比例和平均定位清除时间 $T/N$。和第三问相比，每个已发现频道的记录多保存两组位置：收到过信号的位置 $x^+$ 和没有收到信号的位置 $x^-$，它们在后面裁剪定位区域时要用。源的真实位置、数量、类型、发射方向和接收半径都不进入决策，算法只能通过检测反馈逐步了解它们。

\subsubsection{模型的建立}\label{sec:q4-2}

\textbf{（1）决策内容与目标。}和第三问一样，要决定的是机器狗每一步去哪里、测哪个频道、什么时候清除，目标是在清除全部源的前提下让总耗时 $T$ 尽量小。本问的模型仍分三层：先用固定的检测站保证任何源都会被发现；再对已发现的源用定位区域描述它可能在哪里，并规定哪些反馈可以用来缩小区域；最后把“去哪个站”和“处理哪个源”合成一个路线问题，每走一步重新规划。下面依次建立这三层，重点放在因定向源而改变的地方。

\textbf{（2）能发现定向源的覆盖条件。}第三问只要求每个可能的源位置 $g$ 附近 $1000$ m 内有一个检测站，源就一定会在那里被听到。对定向源，站在 $1000$ m 内还不够，它还得落在发射范围里；而发射方向是未知的，所以要求变成：不管 $u$ 朝哪边，$g$ 附近 $1000$ m 内总得有一个站落在半平面 $u\cdot(x-g)\ge0$ 里。一个简单的充分条件是

\begin{equation}\label{eq:q4-2}
g\in\operatorname{conv}\{s\in P:\ \|s-g\|\le1000\},\qquad\text{对每个 } g\in\Omega ,
\end{equation}

其中 $P$ 是检测站集合，$\operatorname{conv}$ 表示凸包。理由见图\ref{fig:q4-cover}(b)：发射范围的边界是一条过 $g$ 的直线，如果 $g$ 附近的站全在这条直线的另一侧，那么它们的凸包也全在另一侧，就不可能把 $g$ 包在里面。所以只要式\eqref{eq:q4-2} 成立，过 $g$ 的任何直线都至少有一个附近的站落在 $u$ 那一侧，它离 $g$ 不到 $1000$ m，又在发射范围内，按式\eqref{eq:q4-1} 一定能收到信号。全向源只要附近有站就行，自然也满足。

按这个条件布站时，边界朝外的源决定了必须有站在场外。本文采用原点、内环 $8$ 站（半径 $998$ m）和外环 $12$ 站（半径 $1865$ m）共 $21$ 个检测站，见图\ref{fig:q4-cover}(a)。外环放在场外是为了让圆周上的点也落进凸包：相邻两个外环站相隔 $30^\circ$，它们连线的中点到原点的距离是 $1865\cos15^\circ\approx1801.5$ m，刚好把 $1800$ m 的圆周包在里面，半径再小一点，圆周上的点就会露到凸包外面。内环半径取 $998$ m，比 $1000$ m 略小，是让原点附近的点能同时用上原点和整个内环。条件\eqref{eq:q4-2} 要对场地里每一个点都成立，程序把场地切成小方格逐格检验：对一个方格，先找出到方格里每一点都不超过 $1000$ m 的站，若它们的凸包把方格四个角都包住，这个方格就通过；不通过的方格一分为四继续检验。这套布局共检验通过 $4228$ 个方格，检验用的是保守的浮点比较，不依赖抽样。

有了覆盖条件，发现规则就很简单：每到一个站，把所有还没发现的频道都测一遍。第三问里靠“可以证明无信号”省掉的检测，本问对未知频道一律不省，因为无信号在这里推不出任何几何结论。只要 $21$ 个站都扫过，任何还没出现过的频道就一定没有源；已经发现 $16$ 个频道后，剩下的站和频道也不必再扫。

\begin{center}
\begin{minipage}{\linewidth}
\centering
\QFourCoverageFigure
\captionof{figure}{发现覆盖的布局与凸包条件。(a) 原点、内环 $8$ 站与外环 $12$ 站共 $21$ 个检测站，红色虚线圆是候选源位置 $g$ 的 $1000$ m 接收范围，绿色区域是范围内各站的凸包。(b) 放大 $g$ 附近：$g$ 在凸包内，所以不管发射方向 $u$ 朝哪边，总有一个范围内的站落在发射范围里。}
\label{fig:q4-cover}
\end{minipage}
\end{center}

\textbf{（3）有信号时的更新。}某个频道第一次在站 $s$ 收到示向度 $\theta$ 时，定位区域取楔形 $W(s,\theta)$、以 $s$ 为圆心的 $1500$ m 圆和场地三者的交，多边形外包的做法与第二问相同。以后每收到一次示向度，就再交上一个楔形；收到近场反馈就直接清除。这一部分对两类源完全一样：收到信号时式\eqref{eq:q4-1} 的两个条件都已经满足，方向条件不再提供新的信息。

\textbf{（4）成对检测点：让无信号也能用。}既然单个无信号点没有用，本文改为一次布两个检测点，让“两个点都无信号”能推出结论。设上一次收到信号的位置为 $s$，示向方向的单位向量为 $e$，$v$ 是与 $e$ 垂直的单位向量，取

\begin{equation}\label{eq:q4-3}
q^{\pm}=s+t\,e\pm h\,v,\qquad
\frac ht>\tan\delta,\qquad
\Bigl(\frac ht\Bigr)^2+2\,\frac ht\tan\delta<1 ,
\end{equation}

也就是沿示向方向前进 $t$，再向两侧各错开 $h$，见图\ref{fig:q4-pair}。两个条件分别保证：两点横跨整个 $\pm\delta$ 的楔形；只要源沿 $e$ 的坐标不小于 $t$，两点到源的距离都比 $s$ 到源的距离小。于是可以断言：\textbf{若 $q^+$、$q^-$ 都无信号，则源沿 $e$ 的坐标小于 $t$。}理由是，假如源在 $t$ 之外，两点都比 $s$ 更靠近源，而 $s$ 已经收到过信号，所以两点都在接收范围内；另一方面源在两点张成的锥形里，任何一个包含 $s$ 的半平面都会至少包含 $q^+$、$q^-$ 之一（否则连接 $s$ 和源的线段会穿过 $q^+q^-$，而线段两端都在半平面里，中间的点不可能在外面），所以至少有一个点也在发射范围内，两点不可能同时无信号。这个推断对全向源和定向源都成立，不需要知道 $u$。两点都无信号时，用半平面 $e\cdot(p-s)\le t$ 裁剪定位区域；有一个点收到示向度，就按（3）交上新的楔形，并以该点作为新的 $s$。

程序里 $t$ 取当前区域沿 $e$ 的最远距离的 $30\%$；如果区域沿 $e$ 的最近距离已经超过这个值，就把 $t$ 提到最近距离的 $98\%$（但不超过最远距离的 $95\%$），免得测在明知没有源的地方；$h$ 取最远距离的 $2.5\%$，并保证 $h/t$ 至少是 $1.05\tan\delta$。这时 $h/t\approx0.08$，式\eqref{eq:q4-3} 的两个条件都满足。先测离当前位置近的那个点，收到信号就不必测另一个。每个源最多做 $12$ 轮这样的检测。

\begin{center}
\begin{minipage}{\linewidth}
\centering
\QFourPairFigure
\captionof{figure}{成对检测点及其推断。两点沿示向方向前进 $t$、向两侧错开 $h$；若源在 $t$ 之外，两点都在接收范围内，且至少一点在发射范围内，所以两点同时无信号只可能是源在 $t$ 之内。}
\label{fig:q4-pair}
\end{minipage}
\end{center}

\textbf{（5）用有信号点和无信号点一起裁剪。}无信号点还有一种用法。接收半径至少 $1000$ m，而源到每个有信号点的距离都不超过 $R_f$，所以由定位区域能算出 $R_f$ 的一个下界 $R_{\min}$：取 $1000$ m 和“各有信号点到区域包围圆圆心的距离减去包围圆半径”中的最大者。如果某个无信号点 $x^-$ 到定位区域里每一点的距离都小于 $R_{\min}$，那么它无信号就不可能是因为距离，只能是因为它在发射范围之外（顺便也说明这个源是定向的）。这时对源的位置 $p$ 和方向 $u$ 同时有约束：对每个有信号点 $x^+$ 有 $u\cdot(x^+-p)\ge0$，对 $x^-$ 有 $u\cdot(x^--p)<0$。两式相减得 $u\cdot(x^+-x^-)>0$，即 $u$ 与向量 $x^+-x^-$ 的夹角小于 $90^\circ$，可能的发射方向被限制在半个圆周内；对其中每一个方向 $u$，$p$ 又被夹在过 $x^-$、$x^+$ 且垂直于 $u$ 的两条直线之间，见图\ref{fig:q4-negative}。程序把允许的方向分成 $16$ 个扇区，每个扇区用它的中心方向做裁剪，再按扇区宽度向外放宽一点，最后把各扇区的结果并起来取凸包。这样得到的仍是包住真实源的外包多边形，只是比原来小。裁剪只在 $x^-$ 确实落在 $R_{\min}$ 之内时进行，条件不满足就保持原区域不动，宁可保守，也不冒把真实源裁掉的风险。

\begin{center}
\begin{minipage}{\linewidth}
\centering
\QFourNegativeFigure
\captionof{figure}{定向源无信号反馈的裁剪。无信号点 $x^-$ 离定位区域每一点都不到接收半径下界时，它只能在发射范围之外；由此限定发射方向的范围，并对每个允许方向把源夹在两条垂线之间，条带以外的部分被裁掉。}
\label{fig:q4-negative}
\end{minipage}
\end{center}

\textbf{（6）清除判据与光学覆盖。}清除的判据与第三问相同：定位区域包围圆半径不超过 $r_c$ 时，走到离区域每个顶点都不超过 $r_c$ 的位置就一定能清除，程序在这些位置里取离当前位置最近的一个；本问取 $r_c=19.5$ m，比第三问的 $19.9$ m 再多留一点余量。收到近场反馈的源，走到离该点 $14.5$ m 以内清除即可，因为源离该点不到 $5$ m。定位区域已经很窄、但沿长轴还较长时，不必再测：只要区域能装进一个半宽小于 $19.5$ m 的长条，就沿长轴摆放半径 $19.5$ m 的圆盘把长条盖住，从离当前位置近的一端起逐个圆心尝试光学清除，每次失败按题意计 $3$ s，成功即止；圆盘不超过 $12$ 个时才采用这种做法。真实的源必在某个圆盘内，所以这个过程一定会成功。任何频道只有收到真实的清除成功反馈才登记为已清除。

\textbf{（7）任务的安排与共享测量。}每一步的任务集由两部分组成：还没去过的检测站，以及已发现但未清除的源，源的位置用当前定位区域的包围圆圆心代表，有近场反馈的用近场点。任务最多 $21+16=37$ 个，本问不再用第三问的状态压缩搜索，而用更省时的办法：从离当前位置最近的 $8$ 个任务分别出发做最近邻排序，各自再用 2-opt 改进，取总路程最短的一条，然后只执行它的第一个任务。执行完一个站的扫描、一轮成对检测或一次清除之后，任务集和定位区域都变了，就重新排序。这样做没有最优性保证，但每次规划只需几毫秒，方便后面在一步之内反复调用。

到站扫描未知频道以后，机器狗已经站在那里，对附近的已知源顺便再测一次往往很划算，这就是第三问的共享测量。本问的条件是：该源的包围圆半径超过 $19.5$ m，圆心离当前位置不超过 $600$ m 加上半径，当前位置离这个频道以前的检测点至少 $50$ m，并且按“源在区域中心和长轴两端附近”三个代表位置估算，一次测向预计能把包围圆半径缩小 $10$ m 以上。估算时按各代表位置与已有的有信号点、无信号点的几何关系，给每个位置一个能被收到的可能性，只把可能收到的部分计入缩小量。清除一个源之后、以及成对检测的每一步之后，也用同样的规则检查一遍。

\textbf{（8）处理一个源时的局部前瞻。}轮到处理某个源、而它的区域还不能直接清除时，标准动作是（4）里的一轮成对检测。但成对检测的 $t$、$h$ 是固定比例，不一定是眼下最省时间的走法。本文在执行之前先比较十三个候选动作：标准的一轮成对检测；改变前进比例和横向比例的六种成对检测，$t$、$h$ 与最远距离之比分别取 $(0.3,0.08)$、$(0.5,0.035)$、$(0.65,0.06)$、$(0.8,0.04)$、$(0.4,0.15)$、$(0.65,0.15)$；以及在区域沿 $e$ 的 $1/4$、$1/2$、$3/4$ 处向两侧各错开一点的六个单点检测。比较的办法是模拟：在定位区域里抽 $24$ 个和已有观测不矛盾的假想源，每个假想源都要经得起检验，也就是过去每次有信号、无信号的反馈在它身上都能重现；对每个假想源分别执行候选动作，然后按标准规则一直做到清除该源，再加上走到下一个任务的时间，取 $24$ 次的平均耗时作为该动作的分数。抽假想源时，位置在区域内均匀抽取，类型、发射方向和接收半径按均匀分布考虑，并按它们与历史反馈的相符程度加权。只有当最好的候选比标准动作平均省 $2$ s 以上时才换用它；单点检测全局最多 $40$ 次，这种比较每局最多做 $40$ 次、总计算时间不超过 $120$ s，超出预算就直接执行标准动作。模拟只用来挑动作，真实的定位区域仍然只用真实反馈更新。

\textbf{（9）途中停下检测。}前面的规则都是“到了任务点再测”。实际上从当前位置到下一个任务点往往要走几百米，途中经过一些已知源的附近，停下来对它们测一次并不增加路程，只多花检测和换频的时间；测到的示向度会缩小这些源的区域，之后处理它们时就能少走路、少做几轮成对检测。图\ref{fig:q4-transit} 画出了这个想法。具体规则是：设当前位置为 $p$，下一个任务的入口点为 $r$（检测站就是站的位置；源任务取它的第一个物理到达点，即清除点、光学覆盖的端点或成对检测中较近的那个点），路段不短于 $200$ m 时，考虑在 $1/4$、$1/2$、$3/4$ 三个位置停下。每个停点上可测的频道是这样的已知源：包围圆半径至少 $35$ m，圆心到停点不超过 $800$ m 加半径，停点离该频道以前的检测点至少 $60$ m，而且按（7）的估算一次测向能把半径缩小 $20$ m 以上。停点测到示向度就交上新楔形，无信号只记入 $x^-$ 列表，不套用成对检测的截断结论。

难的是判断停下来到底划不划算：区域缩小并不直接等于省时间，检测本身也要花时间。本文没有为此再设计规则，而是用数据来定：离线生成大量模拟场景，在各种状态下分别记录“不停、直接按标准规则做完全部任务”和“停下检测再做完”的总耗时，二者之差

\begin{equation}\label{eq:q4-4}
Y=T_{\text{不停}}-T_{\text{停}}
\end{equation}

就是停这一下的实际收益。用停点的 $57$ 个可观测特征（路段长度、停点位置、各候选频道的区域大小、预计缩小量、已有正负观测个数、剩余站数、已知源数等）去拟合 $Y$，得到一个由 $160$ 棵回归树组成的模型。在线运行时，对三个候选停点分别算特征、预测 $Y$，取预测值最大的停点，预测值大于零才停。训练样本来自 $2380$ 个模拟场景共 $28359$ 个候选状态，训练时按源数的倒数加权，使目标和题目的 $T/N$ 指标一致；特征里没有任何真实源的位置、类型和方向。为了限制风险，每局最多停 $16$ 次，同一个任务不能连续两次被插入停点，而且在最初四个检测站（含原点）扫完之前不启用这项规则，让搜索初期完全按（2）到（8）的规则走。后一条是开发中的教训：学习得到的评分改变早期路线后，有时会推迟发现最后几个频道；对源数正好是 $16$ 的场景，这意味着失去“发现 $16$ 个就可以停止扫站”的机会，损失可以很大。

\begin{center}
\begin{minipage}{\linewidth}
\centering
\QFourTransitFigure
\captionof{figure}{去下一任务途中的停测。停在线段上不增加路程，只增加检测和换频时间；停不停、停在哪一点、测哪些频道由回归树预测的总耗时节省决定。}
\label{fig:q4-transit}
\end{minipage}
\end{center}

\textbf{（10）终止条件。}任务在两种情况下结束：已经真实清除了 $16$ 个不同频道的源；或者 $21$ 个站都扫完（未知频道在每个站都测过）、已发现的源全部清除。第一种由源数上限保证不会有遗漏，第二种由覆盖条件\eqref{eq:q4-2} 保证。只发现了 $16$ 个源但还没清完，或者只是站扫完了而已知源还在，都不能结束。

\subsubsection{模型的求解}\label{sec:q4-3}

模型里没有一个能一次解出来的优化问题，求解就是按上面的规则把每一步走出来：每执行一步，用真实反馈更新定位区域和任务集，再决定下一步。流程见图\ref{fig:q4-flow}，步骤归纳为算法\ref{alg:q4-search}。整套流程每一步都只用四种真实反馈：示向度、近场、无信号和清除成功与否。几何部分保证不漏源、不裁掉真实源、判定能清除时一定能清除；局部前瞻和回归树评分只在允许的动作里挑一个，不改变这些保证，它们算错了最多是多花时间，不会漏掉源。程序用纯 Python 实现，不依赖第三方库，回归树的权重随程序附带，在线运行不需要再训练。

\begin{center}
\begin{minipage}{\linewidth}
\centering
\QFourFlowFigure
\captionof{figure}{问题四的求解流程。每执行一步都重新规划；前四个检测站扫完之前不启用途中停测。}
\label{fig:q4-flow}
\end{minipage}
\end{center}

\par\medskip
\noindent\begin{minipage}{\linewidth}
\small
\refstepcounter{algorithm}\label{alg:q4-search}
\textbf{算法\thealgorithm：混合定向源的滚动搜索与清除}\par\smallskip
\noindent\textbf{输入：}目标圆域 $\Omega$、$20$ 个频道、源数范围 $10$ 到 $16$、$21$ 个检测站、设备动作费用
\begin{enumerate}
\setlength{\itemsep}{2pt}
\setlength{\parskip}{0pt}
\item 在原点扫描全部 $20$ 个频道；收到示向度的频道建立定位区域，无信号的位置记入该频道的 $x^-$ 列表；
\item\label{step:q4-check} 若已真实清除 $16$ 个频道，或 $21$ 站扫完且已发现的源全部清除，结束；
\item 把未访问的站和已发现未清除的源组成任务集，按最近邻加 2-opt 排序，取首任务；
\item 若已扫完至少 $4$ 个站：对去首任务途中的三个候选停点算特征并用回归树预测节省，预测值最大且大于零时前往该停点，检测选定的频道并更新区域，回到步骤\ref{step:q4-check}；
\item 首任务是检测站：前往该站，扫描全部未知频道，对附近满足（7）条件的已知源补测，回到步骤\ref{step:q4-check}；
\item 首任务是源：若包围圆半径不超过 $19.5$ m 或能用不超过 $12$ 个圆盘做光学覆盖，前往清除并按真实反馈登记；否则由局部前瞻在标准成对检测、其他比例的成对检测和单点检测中选一个执行，收到示向度就交上新楔形，两点都无信号就用 $e\cdot(p-s)\le t$ 截断区域，并按（5）检查能否用无信号点再裁剪；回到步骤\ref{step:q4-check}。
\end{enumerate}
\end{minipage}
\par\medskip

\subsubsection{结果分析}\label{sec:q4-4}

\textbf{（1）实验设置。}本问没有官方演练日志可用，实验在按题目附件 2 的接口说明和演练程序的生成规则复原的本地模拟器上进行：源在半径 $1770$ m 的圆内按面积均匀分布，数量在 $10$ 到 $16$ 之间，频道从 $20$ 个中不放回抽取，定向源的个数在 $1$ 到 $N$ 之间均匀抽取（允许全部为定向），发射方向均匀随机，接收半径在 $1000$ 到 $1500$ m 之间均匀抽取；测向误差是在 $150$ m 网格上平滑变化的空间误差场，先量化到 $0.01^\circ$ 再限制在 $\pm1^\circ$ 内，同一位置、同一频道的误差固定；计时按题目规则：移动 $5$ m/s，检测 $5$ s，换频 $1$ s，光学定位 $3$ s，清除成功再加 $2$ s。它不是官方程序本身，也不能代表正式测试的分布，只用来在相同场景上比较不同策略。

比较三种策略：\textbf{基础方案}，只用第\ref{sec:q4-2}节（2）到（7）的规则，不做局部前瞻，也不在途中停测；\textbf{本文方案}，即算法\ref{alg:q4-search}，在基础方案上加局部前瞻和途中停测；\textbf{完整版}，在本文方案之外再用另一个回归树给任务重新排序，并用模拟得到的后验估计代替（7）中的共享测量判据。完整版是本文方案的前身，这里一并列出，用来说明去掉那两个模块之后效果是否变化。场景分三组：随机主测试 $1000$ 个；压力测试两组各 $70$ 个，一组全部是定向源，另一组所有源都贴着边界朝外发射且接收半径取下限 $1000$ m，两组按源数 $10$ 到 $16$ 各 $10$ 个。三种策略在每个场景上都从同一个初始状态出发，共 $3420$ 次运行，场景在冻结策略之后才生成，没有用来调参数。评价指标是每个场景的总耗时除以该场景的源数，再对场景等权平均；总耗时包括清除最后一个源之后为确认没有遗漏而做的扫描。两种策略的差用 $50000$ 次同场景配对自助法给出 $95\%$ 区间。

\textbf{（2）主要结果。}$3420$ 次运行全部完整清除，共 $14806$ 个源实例，没有失败或超时的场景。三种策略的平均定位清除时间见表\ref{tab:q4-main}。

\Needspace{40mm}
\begin{center}
\begin{minipage}{\linewidth}
\centering\small
\captionof{table}{三种策略的平均定位清除时间与计算时间}
\label{tab:q4-main}
\vspace{4pt}
\begin{tabular}{@{}lrrrrr@{}}
\toprule
策略 & 主测试 / s/源 & 比基础方案节省 & 全定向 / s/源 & 边界朝外 / s/源 & CPU / s/局 \\
\midrule
基础方案 & $454.41$ & --- & $480.05$ & $567.80$ & $0.025$ \\
本文方案 & $449.89$ & $0.996\%$ {[}0.79, 1.20{]} & $470.26$ & $515.43$ & $0.332$ \\
完整版 & $449.55$ & $1.069\%$ {[}0.83, 1.32{]} & $470.03$ & $514.73$ & $0.447$ \\
\bottomrule
\end{tabular}
\end{minipage}
\end{center}

主测试中本文方案比基础方案平均省 $1.0\%$，区间不含零；全定向一组省 $2.0\%$，边界朝外一组省 $9.2\%$，区间为 {[}8.3, 10.1{]}\%。后一组的收益最大，是因为这类源的定位区域又长又窄，基础方案要靠多次光学尝试才能清除，而途中的补测把区域缩小之后光学尝试少了很多。本文方案与完整版之间，主测试的配对差为每源 $0.33$ s，区间 {[}$-0.43$, $1.10${]} 含零；$1000$ 个场景中本文方案更快 $483$ 个、更慢 $449$ 个、相同 $68$ 个，两组压力测试的差也在 $1$ s 以内。也就是说，去掉任务重排评分和后验共享评分之后，平均效果没有可辨别的变化，而每局的计算时间少了约四分之一。这是本文采用简化后方案的依据。

\textbf{（3）时间省在哪里。}表\ref{tab:q4-components} 把主测试的平均耗时按动作拆开。本文方案每局平均少走 $369$ m，光学未命中从 $9.7$ 次降到 $7.1$ 次，这两项省下的时间超过了多做的 $4.4$ 次检测和 $2.7$ 次换频所花的时间。平均每局途中停测 $8.5$ 次、共检测 $11.6$ 次；局部前瞻平均调用 $11.2$ 次，其中 $2.0$ 次换用了非标准动作。检测和换频的总次数只略有增加，说明途中补测替代了一部分本来后面也要做的检测。

\Needspace{36mm}
\begin{center}
\begin{minipage}{\linewidth}
\centering\small
\captionof{table}{主测试中每源平均耗时的分解（单位：s）}
\label{tab:q4-components}
\vspace{4pt}
\begin{tabular}{@{}lrrrrrr@{}}
\toprule
策略 & 移动 & 检测 & 换频 & 光学定位 & 清除 & 合计 \\
\midrule
基础方案 & $330.99$ & $97.66$ & $18.51$ & $5.25$ & $2.00$ & $454.41$ \\
本文方案 & $325.16$ & $99.37$ & $18.72$ & $4.64$ & $2.00$ & $449.89$ \\
\bottomrule
\end{tabular}
\end{minipage}
\end{center}

\textbf{（4）按源数分组。}表\ref{tab:q4-count} 按真实源数（只在事后分组时使用）列出主测试的结果。源越多，每源平均时间越短，因为扫完 $21$ 个站的固定开销被更多的源分摊；源数为 $16$ 的场景平均只访问了 $16.1$ 个站，其余场景都访问了全部 $21$ 个，这是终止条件里“发现 $16$ 个就停止扫站”在起作用。本文方案在七个分组里都比基础方案快，源数为 $16$ 的一组节省最少，因为这些场景本来就常常提前结束。

\Needspace{40mm}
\begin{center}
\begin{minipage}{\linewidth}
\centering\small
\captionof{table}{主测试按源数分组的平均定位清除时间（单位：s/源）}
\label{tab:q4-count}
\vspace{4pt}
\begin{tabular}{@{}lrrrrrrr@{}}
\toprule
源数 & 10 & 11 & 12 & 13 & 14 & 15 & 16 \\
\midrule
场景数 & 151 & 141 & 142 & 139 & 136 & 146 & 145 \\
基础方案 & $588.84$ & $536.63$ & $484.02$ & $448.94$ & $414.16$ & $383.13$ & $320.25$ \\
本文方案 & $582.81$ & $531.15$ & $479.93$ & $443.84$ & $410.05$ & $378.45$ & $318.10$ \\
\bottomrule
\end{tabular}
\end{minipage}
\end{center}

\textbf{（5）个例的退步。}平均改善不等于每个场景都改善。相对基础方案，本文方案最差的一个场景慢 $65.8\%$，慢 $3.4\%$ 以上的场景占 $5\%$。相对完整版最差的场景是编号 main-0066 的一个 $16$ 源场景：完整版在 $2729$ s、访问 $10$ 个站后就发现了全部 $16$ 个频道，随即停止扫站，总耗时 $4071$ s；本文方案早期路线不同，直到 $5376$ s、访问 $20$ 个站后才发现最后一个频道，总耗时 $5678$ s。退步最大的五个场景源数都是 $16$，原因都是这一条：发现最后几个频道的先后，决定了能不能省掉剩下的站。前四站不启用途中停测的规则减轻了这个问题，但没有消除它。

\textbf{（6）计算开销与核验。}表\ref{tab:q4-main} 最后一列是单进程复跑 $50$ 个场景得到的每局 CPU 时间，本文方案为 $0.33$ s，比完整版少 $25.7\%$，局部前瞻的 $120$ s 时间预算在全部运行中一次也没有触发。$3420$ 次会话共 $951142$ 个动作全部按记录重放核对，反馈和计时逐条一致；运行前后核对了源码、模型和配置的哈希值；每个场景都审计了退出条件：源数少于 $16$ 时，所有未发现频道在 $21$ 个站都测过，源数为 $16$ 时有 $16$ 次真实的清除成功。策略只接收检测和清除反馈，不接收真实位置、类型、方向、半径和剩余源数。

\textbf{（7）适用条件。}以上结果只对本文冻结的参数和本地模拟器的场景生成方式成立。途中停测的回归树是在另一套本地模拟场景上训练的，如果官方演练的源分布、误差特性与此差别较大，它的收益可能变小，但几何部分的保证不受影响：无论评分是否准确，$21$ 站扫完就不会漏源，判定能清除时一定能清除。官方演练和正式测试的清除比例与平均定位清除时间，要以实际日志为准。

\subsubsection{问题四的结论}\label{sec:q4-5}

把本节的做法归纳成对题目的回答：

\begin{itemize}
\item \textbf{检测策略。}发现阶段用原点、半径 $998$ m 的内环 $8$ 站和半径 $1865$ m 的外环 $12$ 站共 $21$ 个检测站，每到一站把所有未发现的频道测一遍；任何位置的源，不论全向还是定向、朝哪个方向发射，都会在其中某个站被收到，贴着边界朝外发射的源由场外的外环站负责。定位阶段对每个已发现的源沿示向方向成对布置检测点：收到示向度就交上新楔形，两点都无信号就把定位区域截到前进距离之内；确定在接收范围内却无信号的点，和有信号点一起用来限制发射方向和源的位置。途中经过已知源附近时，由回归树判断是否停下补测。
\item \textbf{清除策略。}定位区域的包围圆半径不超过 $19.5$ m，或区域能用不超过 $12$ 个半径 $19.5$ m 的圆盘盖住时，就前往清除，光学清除不受发射方向影响；只有收到真实的成功反馈才算清除。任务顺序由最近邻加 2-opt 决定，每执行一步重新规划；$16$ 个频道全部清除，或 $21$ 站扫完且已发现的源全部清除时结束。
\end{itemize}

这套策略能保证不漏源、不裁掉真实源、判定能清除时一定能清除。它省了多少时间，以上面的本地演练结果为依据；官方演练和正式测试的成绩要以实际日志为准。
